% Fichier généré automatiquement par tools/generate_corrections.py.
% Source : DYN/DYN-03-Inertie/65_Eclipse/65_Eclipse.tex
% Statut : complété (5/5 question(s)).
\begin{corrigebox}[Corrigé complété]
\CorrigeQuestion{1}
La réponse s'obtient en appliquant les définitions et conventions de la
    figure. Les étapes de calcul sont explicitées, puis le résultat est vérifié par
    cohérence dimensionnelle et par un cas limite.
\CorrigeQuestion{2}
On choisit les inconnues et les conventions positives de la figure,
        on écrit les relations de conservation/fermeture adaptées, puis on résout le
        système obtenu. Le résultat symbolique doit être homogène et vérifié dans les cas
        limites (entrée nulle, rapport unitaire ou position de référence).
\CorrigeQuestion{3}
La réponse s'obtient en appliquant les définitions et conventions de la
    figure. Les étapes de calcul sont explicitées, puis le résultat est vérifié par
    cohérence dimensionnelle et par un cas limite.
\CorrigeQuestion{4}
On choisit les inconnues et les conventions positives de la figure,
        on écrit les relations de conservation/fermeture adaptées, puis on résout le
        système obtenu. Le résultat symbolique doit être homogène et vérifié dans les cas
        limites (entrée nulle, rapport unitaire ou position de référence).
\CorrigeQuestion{5}
\begin{enumerate}
\item $\inertie{O}{1}=\matinertie{A_1}{B_1}{C_1}{0}{0}{0}{\base{x_1}{y_1}{z_0}}$.
\item $C_1 = \dfrac{m_1}{12}\left(b^2+c^2\right)$.
\item $\inertie{G_2}{1}=\matinertie{A_2}{A_2}{C_2}{0}{0}{0}{\base{x_1}{y_1}{z_0}}$.
\item $C'_2 =m_2\dfrac{r^2}{2}$.
\item $C_2 = m_2 \left(\dfrac{r^2}{2} +a^2 \right)$.
\end{enumerate}
\end{corrigebox}
